\section{A direct spectral estimator}
\label{sec:algorithm}
The inverse estimates concern the true experiment, but their proof
does not require a global search over latent factors. We give an
explicit estimator for the multiset alone. Computation of the exact
residual confidence region in the next section is a separate problem.

Given matrices $\widehat P_j$, form $\cN_{\widehat P}$ and
$b_{\widehat P}$ by \eqref{eq:normalization}. If the former has full
column rank, let $\widehat a$ be its least-squares solution and put
$\widehat q(z)=z^d+\sum_{r<d}\widehat a_rz^r$.
For $d=1$, replace $-\widehat a_0$ by its projection onto $[L,D]$;
write the resulting polynomial as $z-\widehat h$.
Interpolate the matrices $\widehat q(T_j)\widehat P_j$ at the first
$d+1$ clocks to obtain $\widehat K(z)$.

Take the leading $k$ left and right singular subspaces of
$\widehat P_1$, with orthonormal bases $\widehat Q_1,\widehat Q_2$.
Form the reduced polynomial
\[
 \widehat M(z)=\widehat Q_1^{\mathsf T}\widehat K(z)\widehat Q_2.
\]
Compute its $kd$ roots, counted with algebraic multiplicity, project
the real parts onto $[L,D]$, and sort them. This is the estimator
$\widehat\Lambda$. If the normalization matrix lacks full column
rank, a sampled $\widehat q$ is nonpositive, $\widehat P_1$ has rank
less than $k$, or the leading coefficient of $\widehat M$ is singular,
return a fixed multiset in $[L,D]$ instead. This convention defines the
estimator everywhere. Numerical root tolerances can be added to the
bounds below; the statements themselves refer to exact arithmetic.

\begin{theorem}[Constructive recovery]\label{thm:algorithm}
Put $e=\max_j\|\widehat P_j-P_j\|_F$.
At every polynomial truth with $\kappa>0$ and $\tau(P)>0$, the
estimator just defined satisfies
\begin{equation}\label{eq:algorithm-poly}
 d_\infty(\widehat\Lambda,\Lambda)
 \le C\min\{1,[e(\kappa^{-1}+\tau(P)^{-1})]^{1/d}\}.
\end{equation}
The within-component simple-root hypothesis of Theorem~\ref{thm:main}
gives the corresponding linear bound. In the affine experiment the
stronger observable estimate holds:
\begin{equation}\label{eq:algorithm-affine}
 d_\infty(\widehat\xi,\xi)
                    \le C\min\{1,e/\eta(P)\},\qquad P\in\mathcal R.
\end{equation}
The procedure uses a least-squares solve, a singular value decomposition,
polynomial interpolation, and the roots of a regular matrix polynomial
of size $k$ and degree $d$. It does not optimize over $U,V$ or enumerate
a parameter net. No conditioning floor is an input to the estimator.
\end{theorem}
\begin{proof}
The normalization operator and its right-hand side depend linearly on
the observed entries. Therefore
$\|\cN_{\widehat P}-\cN_P\|\le Ce$, and the residual of the true
coefficient vector in the perturbed equation is at most $Ce$.
If $e\le c\tau$, the least singular value is at least $\tau/2$.
Least-squares optimality followed by the triangle inequality gives
\begin{equation}\label{eq:algorithm-q}
 \|\widehat a-a\|\le Ce/\tau.
\end{equation}
Projection of $\widehat h$ cannot increase its distance from $h$.
For sufficiently small $e/\tau$, all sampled denominator values
remain positive.

We record the subspace estimate used in both parts of the proof.
Let $Q_1,Q_2$ span the true column and row spaces of $P_1$. The
best rank-$k$ approximation property implies
\[
 \|(I-\widehat Q_1\widehat Q_1^{\mathsf T})P_1\|\le2e,
 \qquad
 \|P_1(I-\widehat Q_2\widehat Q_2^{\mathsf T})\|\le2e.
\]
If $2e/\sigma_k(P_1)$ is sufficiently small, the largest principal
angle has sine at most that ratio. In particular both matrices
$\widehat Q_1^{\mathsf T}Q_1$ and
$Q_2^{\mathsf T}\widehat Q_2$ are invertible with inverse norms at
most two. This assertion also follows directly by applying the
preceding residual inequalities to a singular value decomposition of
$P_1$.

For the polynomial bound, $\sigma_k(P_1)\ge c\kappa$.
Consequently $S=\widehat Q_1^{\mathsf T}U$ and
$T=\widehat Q_2^{\mathsf T}V$ are invertible, with
$\|S^{-1}\|\|T^{-1}\|\le C/\kappa$, when $e/\kappa$ is small.
At the interpolation nodes, reduction by these inverses yields
\begin{align*}
 S^{-1}\widehat M(T_j)T^{-\mathsf T}-D_f(T_j)
 &=\widehat q(T_j)S^{-1}\widehat Q_1^{\mathsf T}
            (\widehat P_j-P_j)\widehat Q_2T^{-\mathsf T}\\
 &\quad +(\widehat q(T_j)-q(T_j))\diag(p_b(T_j)).
\end{align*}
As in \eqref{eq:diagonal-error}, interpolation bounds the coefficient
error by $Ce(\kappa^{-1}+\tau^{-1})$. The leading coefficient is
then invertible and Lemma~\ref{lem:roots} applies. The real-part and
interval projections cannot increase the distance to any matched true
real root. Sorting gives the claimed multiset estimate. All failure
conditions are excluded in this small-error regime; bounded range
handles its complement.

For the stronger affine assertion, no comparison with $\kappa$ is
needed. The resolvent formula \eqref{eq:residue-formula} gives
\[
 \sigma_k(P_1)^{-1}
       =\|(T_1-h)M(T_1)^{-1}\|\le C\beta(P).
\]
Thus the two estimated subspaces have the preceding inverse bounds
as soon as $e\beta$ is small. The true pencil, reduced into these
estimated subspaces, has residues
\[
 (Q_2^{\mathsf T}\widehat Q_2)^{-1}R_x
                (\widehat Q_1^{\mathsf T}Q_1)^{-1},
\]
whose total norm is at most $4\beta$.
By \eqref{eq:algorithm-q} and $\tau\asymp\gamma$,
$u=\widehat h-h$ satisfies $|u|\le Ce/\gamma$.
Apply the exact gauge \eqref{eq:gauge} to this true reduced pencil.
The difference between the estimated pencil and that gauge polynomial
is the affine interpolant of
\[
 (T_j-\widehat h)\widehat Q_1^{\mathsf T}
                      (\widehat P_j-P_j)\widehat Q_2,
 \qquad j=1,2,
\]
and has coefficient norm at most $Ce$. The resolvent and root-counting
argument of Theorem~\ref{thm:residue} now gives root matching error
$C\beta e$ about the shifted actions, plus $Ce/\gamma$ for the
shift itself. It applies to complex estimated roots as well. Projection
and the bounded-range convention complete \eqref{eq:algorithm-affine}.
\end{proof}

The theorem is an arithmetic and stability statement for this explicit
procedure. It is not a bit-complexity theorem uniform in dimension,
precision and conditioning. In degree one only a $k\times k$
generalized eigenvalue problem is required. Confidence-region
optimization is not being declared tractable on the strength of this
point-estimation result.
